\subsubsection*{定理 5.4.}
设 $X$ 是 $\mathbb F_q$ 上维数为 $n$ 的有限型分离概形,
$\underline G$ 是 $X$ 上的可构造 $E$ 层. 设对每个 $x\in X$,
$G_x^\vee$ 的自同态 $F_x$ 是 $T$-整的. 则 $\underline H_c^i(X,\underline G)^\vee$
的自同态 $q^iF^\vee$ 和 $q^nF^\vee$ 都是 $T$-整的.

\medskip
\noindent\emph{证明.}
如定理~5.2.2 中所做, 可归结为 $X$ 是光滑曲线且 $\underline G$ 在 $X$ 上扭常值的情形.
设 $j$ 是 $X$ 到将其紧化的光滑射影曲线 $X'$ 中的包含.

由 Poincar\'e 对偶, $\underline H^i(X',j_*\underline G)$ 与
$\underline H^{2-i}(X',j_*\underline G^\vee(1))$ 对偶.
由推论~5.3(a), 定理~5.2.2 适用于 $j_*\underline G^\vee$:
\begin{align*}
\underline H^0(X',j_*\underline G^\vee)
  &\simeq \underline H^2(X',j_*\underline G)^\vee(-1)
  &&\text{是 \(T\)-整的},\\
\underline H^1(X',j_*\underline G^\vee)
  &\simeq \underline H^1(X',j_*\underline G)^\vee(-1)
  &&\text{是 \(T\)-整的},\\
\underline H^2(X',j_*\underline G^\vee)(1)
  &\simeq \underline H^0(X',j_*\underline G)^\vee
  &&\text{是 \(T\)-整的}.
\end{align*}
这就对 $j_*\underline G$ 证明了定理.

由下列正合列定义 $\underline H$
\[
0\longrightarrow j_!\underline G\longrightarrow j_*\underline G
\longrightarrow\underline H\longrightarrow0.
\]
若 $Y$ 是 $\underline H$ 的支撑, 则上同调长正合列给出
\begin{center}
\begin{tikzcd}[row sep=large,column sep=large]
{} & \underline H_c^0(X,\underline G)
  \arrow[r,hook] & \underline H^0(X',j_*\underline G) & ,\\
\underline H^0(Y,\underline H) \arrow[r] &
  \underline H_c^1(X,\underline G) \arrow[r] &
  \underline H^1(X',j_*\underline G) & ,\\
{} & \underline H_c^2(X,\underline G)
  \arrow[r,"\sim"] & \underline H^2(X',j_*\underline G) & ,
\end{tikzcd}
\end{center}
因此 $(X,\underline G)$ 的定理由 $(X',j_*\underline G)$ 的定理得出.
